\documentclass[11pt]{article}
\usepackage[margin=1in]{geometry}
\usepackage{amsmath,amssymb,amsthm}
\usepackage[utf8]{inputenc}
\usepackage{listings}
\usepackage{hyperref}
\usepackage{xcolor}
\usepackage{textcomp}

\definecolor{leanblue}{rgb}{0.2,0.4,0.7}

\lstdefinelanguage{Lean}{
  keywords={theorem,lemma,def,axiom,structure,class,instance,where,by,sorry,simp,exact,intro,apply,have,let,fun,match,if,then,else,import,open,namespace,end,section,variable,inductive,noncomputable,abbrev,deriving},
  keywordstyle=\color{leanblue}\bfseries,
  comment=[l]{--},
  morecomment=[s]{/-}{-/},
  string=[b]",
  basicstyle=\ttfamily\small,
  breaklines=true,
  extendedchars=true,
  inputencoding=utf8,
  literate=
    {≥}{{$\geq$}}1
    {≤}{{$\leq$}}1
    {→}{{$\to$}}1
    {←}{{$\leftarrow$}}1
    {∧}{{$\wedge$}}1
    {∨}{{$\vee$}}1
    {¬}{{$\neg$}}1
    {∀}{{$\forall$}}1
    {∃}{{$\exists$}}1
    {⊆}{{$\subseteq$}}1
    {⊂}{{$\subset$}}1
    {∈}{{$\in$}}1
    {∉}{{$\notin$}}1
    {×}{{$\times$}}1
    {⊗}{{$\otimes$}}1
    {▸}{{$\blacktriangleright$}}1
    {⟨}{{$\langle$}}1
    {⟩}{{$\rangle$}}1
    {λ}{{$\lambda$}}1
    {α}{{$\alpha$}}1
    {β}{{$\beta$}}1
    {σ}{{$\sigma$}}1
    {θ}{{$\theta$}}1
    {ε}{{$\varepsilon$}}1
    {μ}{{$\mu$}}1
    {π}{{$\pi$}}1
    {Σ}{{$\Sigma$}}1
    {Π}{{$\Pi$}}1
    {▶}{{$\triangleright$}}1
    {⊔}{{$\sqcup$}}1
    {⊓}{{$\sqcap$}}1
    {⊥}{{$\bot$}}1
    {⊤}{{$\top$}}1
    {∅}{{$\emptyset$}}1
    {∪}{{$\cup$}}1
    {∩}{{$\cap$}}1
    {≠}{{$\neq$}}1
    {ℕ}{{$\mathbb{N}$}}1
    {₀}{{$_0$}}1
    {₁}{{$_1$}}1
    {₂}{{$_2$}}1
    {|>}{{$|{\!>}$}}2
    {·}{{$\cdot$}}1
    {—}{{---}}1
    {∘}{{$\circ$}}1
    {√}{{$\sqrt{}$}}1
    {∝}{{$\propto$}}1
    {≈}{{$\approx$}}1
    {═}{{=}}1
    {⟹}{{$\Longrightarrow$}}1,
}

\newtheorem{theorem}{Theorem}
\newtheorem{lemma}[theorem]{Lemma}
\newtheorem{definition}[theorem]{Definition}
\newtheorem{axiom_stmt}[theorem]{Axiom}

\title{Formal Verification: BLS Signature Aggregation}
\author{Zach Kelling, Lux Industries}
\date{December 2025}

\begin{document}
\maketitle

\begin{abstract}
We prove the correctness of BLS signature aggregation used in Quasar consensus. The key result is that aggregating two valid BLS signatures on the same message produces a valid aggregate, following from the bilinearity of the pairing function $e: G_1 \times G_2 \to G_T$.
\end{abstract}

\section{Introduction}
BLS signatures are used in Quasar for fast aggregation. A single 48-byte (BLS12-381 $G_1$ compressed) aggregate signature represents arbitrarily many individual signatures. The bilinear pairing property enables this constant-size aggregation.

\section{Definitions}
\begin{definition}[G1, G2, GT]
Abstract group types for BLS12-381.
\end{definition}

\begin{definition}[Signature]
BLS signature: signature in $G_1$, public key in $G_2$, message hash in $G_1$.
\end{definition}


\section{Theorems and Axioms}
\begin{axiom_stmt}[\texttt{bilinear\_left}]
Left bilinearity: $e(a+b, c) = e(a,c) \cdot e(b,c)$.
\end{axiom_stmt}

\begin{axiom_stmt}[\texttt{bilinear\_right}]
Right bilinearity: $e(a, b+c) = e(a,b) \cdot e(a,c)$.
\end{axiom_stmt}

\begin{theorem}[\texttt{aggregate\_two\_sound}]
If $\text{verify}(s_1)$ and $\text{verify}(s_2)$ on the same message, then $e(\text{sig}_1 + \text{sig}_2, g_2) = e(\text{msg}, \text{pk}_1 + \text{pk}_2)$.
\end{theorem}
\begin{proof}
By bilinearity: $e(\text{sig}_1+\text{sig}_2, g_2) = e(\text{sig}_1,g_2) \cdot e(\text{sig}_2,g_2) = e(\text{msg},\text{pk}_1) \cdot e(\text{msg},\text{pk}_2) = e(\text{msg}, \text{pk}_1+\text{pk}_2)$.
\end{proof}

\begin{axiom_stmt}[\texttt{proof\_of\_possession}]
Rogue key resistance: each validator must prove knowledge of their secret key.
\end{axiom_stmt}

\begin{theorem}[\texttt{bls\_quorum\_sufficient}]
With $3f < n$ and $\text{sigs} \geq 2f+1$, we have $\text{sigs} > n/2$.
\end{theorem}

\begin{theorem}[\texttt{quorum\_overlap}]
Two quorums of $\geq 2f+1$ from $n$ nodes with $3f < n$ satisfy $q_1 + q_2 > n$.
\end{theorem}


\section{Lean Source}
\lstinputlisting[language=Lean]{../../formal/lean/Crypto/BLS.lean}

\section{References}
\noindent Source code: \url{https://github.com/luxfi/proofs/blob/main/lean/Crypto/BLS.lean}\\
\noindent White papers: \url{https://papers.lux.network}

\noindent Related white papers:
\begin{itemize}
  \item \texttt{lux-quasar-consensus.tex}
\end{itemize}

\end{document}
